\documentclass[12pt]{beamer}
% \usetheme{umbc4}
\usetheme{moloch}

\usepackage{amsmath,amscd}
\usepackage{fontspec}
\usepackage{unicode-math}
\usepackage{pifont}
\usepackage{xcolor}
\usepackage{colortbl}
\usepackage[dvipsnames]{xcolor}
\usepackage{listings}
\usepackage{minted}
\usepackage{tikz}
\usepackage{hyperref}
\usepackage[normalem]{ulem}

\hypersetup{
    colorlinks=true,
    linkcolor=blue,       % internal document links
    citecolor=cyan,       % bibliography citations
    urlcolor=lightblue    % web URLs
}

\setmonofont{DejaVu Sans Mono}

\setbeamercolor{frametitle}{bg=black, fg=white}

% Ensure other structural elements (bullets, etc.) are visible
\setbeamercolor{structure}{fg=white}

\setmainfont{Libertinus Serif}
\setsansfont{Libertinus Sans}
\setmathfont{STIX Two Math}

% Define colors for syntax highlighting
\definecolor{codegreen}{rgb}{0,0.6,0}
\definecolor{codegray}{rgb}{0.5,0.5,0.5}
\definecolor{codecoffee}{rgb}{0.75,0.6,0.3}
\definecolor{question}{rgb}{1.0,0.75,0.1}
\definecolor{lightblue}{rgb}{0.55,0.75,0.9}
\definecolor{yellow}{rgb}{1.0, 1.0, 0.0}
\definecolor{codecomment}{rgb}{0.97, 0.74, 0.4}
\definecolor{codepy}{rgb}{0.9, 0.85, 0.85} %{#F8BC65}
\definecolor{funcColor}{rgb}{0.47, 0.6, 0.8}
\definecolor{midgray}{rgb}{0.4, 0.4, 0.3}
\definecolor{indexwhite}{rgb}{0.9, 0.9, 0.9}
\definecolor{darkgold}{HTML}{5D3D00}

\setbeamercolor{background canvas}{bg=black}
\setbeamercolor{normal text}{fg=white}

\newcommand{\FrameTitle}[2]{\frametitle{#1 \hfill \small\textnormal{\textcolor{midgray}{#2}}}}
\newcommand{\snl}{\vspace*{6pt}\\}
\newcommand{\nl}{\vspace*{1em}\\}
\newcommand{\vsp}{\vspace*{1em}}

\lstdefinestyle{mypystyle}{
    commentstyle=\color{codecomment},
    keywordstyle=\color{lightblue},
    numberstyle=\scriptsize\color{codegray},
    stringstyle=\color{codecoffee},
    basicstyle=\ttfamily\scriptsize\color{codepy},
    breakatwhitespace=false,
    breaklines=true,
    lineskip=2pt,  % introduces a small gap between lines for better readability
    % captionpos=b,
    keepspaces=true,
    numbers=left,
    numbersep=5pt,
    firstnumber=1,  % Forces reset to 1 for EVERY listing
    showspaces=false,
    showstringspaces=false,
    showtabs=false,
    tabsize=2,
    language=Python,
    backgroundcolor=\color{black},
    emph=[2]{__init__, calculate},
    emphstyle=[2]\color{funcColor},
}
\lstset{style=mypystyle}

\begin{document}
\title{\center{Applied Linear Algebra\\}}

\author{\center{Devi Prasad M \\ Manipal School of Information Sciences \\ Manipal \\\vsp\vsp\color{orange!45}{https://github.com/DeviPrasad/ala2026}}}
\date{\center{Monday, 3 August 2026}}
\maketitle

\begin{frame}\FrameTitle{The Basics}{Introduction}
    \begin{itemize}
        \item[] {\Large\textbf{What is algebra?}}
        \item[]
        \item[] {\Large\textbf{What is linear algebra?}}
        \item[]
        \item[] {\Large\textbf{What is applied linear algebra?}}
    \end{itemize}
\end{frame}

\begin{frame}\FrameTitle{The Basics}{Introduction}
    \begin{itemize}
        \item[] {\Large\textbf{What is algebra?}}
        \begin{itemize}
            \item[] \href{https://en.wikipedia.org/wiki/Algebra}{Algebra - Wikipedia}
        \end{itemize}
    \end{itemize}
    \vspace*{1em}
    \begin{itemize}
        \item[] {\Large\textbf{What is linear algebra?}}
        \begin{itemize}
            \item[] \href{https://en.wikipedia.org/wiki/Linear_algebra}{\color{yellow!40}Linear Algebra - Wikipedia}
        \end{itemize}
    \end{itemize}
\end{frame}

\begin{frame}\FrameTitle{The Basics}{Introduction}
    \begin{itemize}
        \item[] {\Large\textbf{What is applied linear algebra?}}
        \begin{itemize}
            \item[---] \href{https://en.wikipedia.org/wiki/Numerical_linear_algebra}{\color{orange!40}Numerical Linear Algebra - Wikipedia}
            \vspace*{2pt}
            \item[---] \href{https://comp-lin-alg.github.io/}{\color{orange!40}Computational Linear Algebra - GitHub}
        \end{itemize}
    \end{itemize}
\end{frame}

\begin{frame}\FrameTitle{Linear Algebra}{Introduction}
    Consider the set of linear algebraic equations:
    \begin{align*}
        a_{00}x_{0} + a_{01}x_{1} + a_{02}x_{2} + \ldots + a_{0,N-1}x_{N-1} = b_0\\
        a_{10}x_{0} + a_{11}x_{1} + a_{12}x_{2} + \ldots + a_{1,N-1}x_{N-1} = b_1\\
        a_{20}x_{0} + a_{21}x_{1} + a_{22}x_{2} + \ldots + a_{2,N-1}x_{N-1} = b_2\\
        \ldots\\
        a_{M-1,0}x_{0} + a_{M-1,1}x_{1} + \ldots + a_{M-1,N-1}x_{N-1} = b_{M-1}\\
    \end{align*}

    \vspace*{-1em}{\Large\textbf{\color{question}{Quiz:}}}\vspace*{6pt}
    \begin{enumerate}
        \item[1.] How many equations do you see here?
        \item[2.] How many unknowns do we have in this system?
        \item[3.] What is each $a_{ij}$ referred to as?
    \end{enumerate}
\end{frame}

\begin{frame}\FrameTitle{Linear Algebra}{Introduction}
    Consider the set of linear algebraic equations:
    \begin{align*}
        a_{00}x_0 + a_{01}x_1 + a_{02}x_2 + \ldots + a_{0,N-1}x_{N-1} = b_0\\
        a_{10}x_0 + a_{11}x_1 + a_{12}x_2 + \ldots + a_{1,N-1}x_{N-1} = b_1\\
        a_{20}x_0 + a_{21}x_1 + a_{22}x_2 + \ldots + a_{2,N-1}x_{N-1} = b_2\\
        \ldots\\
        a_{M-1,0}x_0 + a_{M-1,1}x_1 + \ldots + a_{M-1,N-1}x_{N-1} = b_{M-1}\\
    \end{align*}

    \vspace*{-1em}{\Large\textbf{\color{question}{Quiz:}}}\vspace*{6pt}
    \begin{enumerate}
        \item[4.] What is each $x_j$?
        \item[5.] What do each $b_m$ represent?
    \end{enumerate}
\end{frame}

\begin{frame}\FrameTitle{Linear Algebra}{Introduction}
    Consider the set of linear algebraic equations:
    \begin{align*}
        a_{00}x_0 + a_{01}x_1 + a_{02}x_2 + \ldots + a_{0,N-1}x_{N-1} = b_0\\
        a_{10}x_0 + a_{11}x_1 + a_{12}x_2 + \ldots + a_{1,N-1}x_{N-1} = b_1\\
        a_{20}x_0 + a_{21}x_1 + a_{22}x_2 + \ldots + a_{2,N-1}x_{N-1} = b_2\\
        \ldots\\
        a_{M-1,0}x_0 + a_{M-1,1}x_1 + \ldots + a_{M-1,N-1}x_{N-1} = b_{M-1}\\
    \end{align*}

    \vspace*{-1em}{\Large\textbf{\color{question}{Quiz:}}}\vspace*{6pt}
    \begin{enumerate}
        \item[6.] \textbf{Why do we care about such a system of equations?}
    \end{enumerate}
\end{frame}

\begin{frame}\FrameTitle{Linear Algebra}{Introduction}
    At the beginning of a semester, {\color{blue!20}{55}} students have signed up for ALA.
    The course is offered in two sections. Because of personal preferences,
    {\color{blue!20}{20\%}} of students in section A switch to section B, while {\color{blue!20}{30\%}} of
    students in section B switch to section A, resulting in a net loss
    of 4 students for section B.

    \vsp
    {\color{question}{How large are the two sections at the beginning of the semester?}}$^\dagger$

    \vfill
    $^\dagger${\color{yellow!20}{\scriptsize{No student dropped ALA or joined the course later.}}}
\end{frame}

\begin{frame}\FrameTitle{Linear Algebra}{Introduction}
    At the beginning of a semester, {\color{blue!20}{55}} students have signed up for ALA.
    The course is offered in two sections. Because of personal preferences,
    {\color{blue!20}{20\%}} of students in section A switch to section B, while {\color{blue!20}{30\%}} of
    students in section B switch to section A, resulting in a {\emph{\uline{net loss
    of {\color{blue!20}{4}} students for section B}}}.

    \vsp
    {\color{question}{How large are the two sections at the beginning of the semester?}}
    \begin{align*}
        a + b & = 55 \tag{1}\\
        b + 0.2a - 0.3b & = \textbf{?} \notag \\
        {\color{yellow}{\vdots}}
    \end{align*}
\end{frame}

\begin{frame}\FrameTitle{Linear Algebra}{Introduction}
    At the beginning of a semester, {\color{blue!20}{55}} students have signed up for ALA.
    The course is offered in two sections. Because of personal preferences,
    {\color{blue!20}{20\%}} of students in section A switch to section B, while {\color{blue!20}{30\%}} of
    students in section B switch to section A, resulting in a net loss
    of 4 students for section B.

    \vsp
    {\color{question}{How large are the two sections at the beginning of the semester?}}
    \begin{align*}
        a + b & = 55 \tag{1}\\
        b + 0.2a - 0.3b & = b-4 \notag \\
        2a - 3b & = -40 \tag{2}\\
        {\color{yellow}{\vdots}}
    \end{align*}
\end{frame}


\begin{frame}[fragile]\FrameTitle{Solving Systems of Linear Equations}{Introduction}
    \vspace*{-2em}
    \begin{align*}
        2x_1 + 8x_2 + 4x_3  &= 2 \\
        2x_1 + 5x_2 + 1x_3  &= 5 \\
        4x_1 + 10x_2 - x_3  &= 1 \\
    \end{align*}
    \vspace*{-1.5em}
    \begin{lstlisting}[language=Python, escapeinside={**@}{@**}]
        import numpy as np
        A = np.array([[2, 8, 4],[2, 5, 1],[4, 10, -1]])
        b = np.array([2, 5, 1])
        x = **@\colorbox{darkgold}{np.linalg.solve}@**(A, b)
        assert np.allclose(np.dot(A, x), b)
        assert np.equal(x, np.array([11, -4, 3])).all()
        assert (x == np.array([11, -4, 3])).all()\end{lstlisting}
    \hfill
\end{frame}

\begin{frame}\FrameTitle{Solving Systems of Linear Equations}{Introduction}
    \vspace*{-2em}
    \begin{align*}
        2x_1 + 8x_2 + 4x_3  &= 2 \\
        2x_1 + 5x_2 + 1x_3  &= 5 \\
        4x_1 + 10x_2 - x_3  &= 1 \\
    \end{align*}

    \vspace*{-1.5em}
    \[
        \underbrace{
            \begin{bmatrix}
                2 & 8 & 4 \\
                2 & 5 & 1 \\
                4 & 10 & -1
            \end{bmatrix}
        }_{\mathbf{A}:\mathbb{R}^{3 \times 3}}
        \,
        \underbrace{
            \begin{bmatrix}
                x_1 \\
                x_2 \\
                x_3
            \end{bmatrix}
        }_{\textbf{x}: \mathbb{R}^{3 \times 1}}
        =
        \underbrace{
            \begin{bmatrix}
                2 \\
                5 \\
                1
            \end{bmatrix}
        }_{\textrm{\textbf{b}}: \mathbb{R}^{3 \times 1}}
    \]
\end{frame}


\begin{frame}[fragile]\FrameTitle{Solving Systems of Linear Equations}{Introduction}
    \vspace*{-2em}
    \begin{align*}
                         x_3 - x_4 - x_5 &= 4 \\
        2x_1 + 4x_2 + 2x_3 + 4x_4 + 2x_5 &= 4 \\
        2x_1 + 4x_2 + 3x_3 + 3x_4 + 3x_5 &= 4 \\
        3x_1 + 6x_2 + 6x_3 + 3x_4 + 6x_5 &= 6 \\
    \end{align*}

    \vspace*{-1.5em}

    \begin{lstlisting}[language=Python, escapeinside={(*@}{@*)}]
        import numpy as np
        A = np.array([[0, 0, 1, -1, -1],
                      [2, 4, 2, 4, 2],
                      [2, 4, 3, 3, 3],
                      [3, 6, 6, 3, 6]])
        b = np.array([4, 4, 4, 6])
        x, _residuals, _rank, _s = (*@\colorbox{darkgold}{np.linalg.lstsq}@*)(A, b)
        assert np.allclose(np.dot(A, x), b)\end{lstlisting}
\end{frame}

\begin{frame}[fragile]\FrameTitle{Solving Systems of Linear Equations}{Introduction}
    \begin{definition}[\color{orange!50}{Underdetermined systems (M < N)}]
        -- There are fewer equations than unknowns. \vspace*{2pt}\\
        -- Infinitely many exact solutions exist. \vspace*{2pt}\\
        -- so the system is underconstrained.
    \end{definition}
    \vspace*{3pt}
    \begin{definition}[\color{orange!50}{Overdetermined systems (M > N)}]
        -- There are more equations than unknowns. \vspace*{2pt}\\
        -- No single vector {\textrm{\textbf{x}}} satisfies $Ax = b$ exactly.  \vspace*{2pt}\\
        -- The vector \textbf{b} simply lies outside the column space of A.
    \end{definition}
\end{frame}
\end{document}
